\documentclass[10pt,a4paper]{amsart}
\usepackage[margin=19mm]{geometry}
\usepackage{amsmath,amssymb,mathtools}
\usepackage[scr=rsfs,cal=euler]{mathalfa}
\usepackage{xfrac}
\usepackage[unicode,hidelinks,bookmarks=false]{hyperref}
\newcommand{\ie}{i.e.\ }
\newcommand{\cf}{cf.\ }
\newcommand{\resp}{respectively\ }
\newcommand{\Subsection}{\subsection}
\providecommand{\nobreakdash}{\nobreak\hbox{-}}
\setlength{\emergencystretch}{2em}
\multlinegap0pt
\usepackage{amssymb}
\usepackage{xcolor}
\usepackage{tikz}
\newtheorem{lemma}{Lemma}
\usepackage{cleveref}
\crefname{lemma}{Lemma}{Lemmas}
\newcommand{\Q}{\mathbb{Q}}
\title{On invariants of representations of Weyl groups associated with the cohomology of toric varieties}
\author{Tao Gong}
\date{Source: arXiv:2602.23631v1 (CC BY 4.0)}
\begin{document}
\maketitle

\noindent\emph{Source and licence.} On invariants of representations of Weyl groups associated with the cohomology of toric varieties. Version arXiv:2602.23631v1 (CC BY 4.0), distributed by arXiv under the Creative Commons Attribution 4.0 International licence, http://creativecommons.org/licenses/by/4.0/, as recorded in the arXiv licence metadata for this exact version and shown on the abstract page https://arxiv.org/abs/2602.23631v1. Full licence notice: \texttt{licenses/arxiv-2602.23631-CC-BY-4.0.txt}.\par\smallskip

\noindent\textit{Editorial scope.} Counted unit: the article's lemma lem:image (source0.tex lines 741--749). The article's complete original proof of it is delivered with it (source0.tex lines 750--778, one \texttt{proof} environment of 29 lines). No mathematics is written, repaired or re-derived: the statement and the proof are the article's own lines, in the article's own notation, and no retained line is substituted or re-flowed.  Retained uncounted as the source-local context the proof needs: lines 342--348; lines 732--733.  Nothing was omitted from any retained span, so the package carries no omission record.  No retained line contains a citation, so the wrapper carries no bibliography and no bibliography entry.  No retained line contains a figure environment, an image-inclusion command or drawing code, so no image is delivered and the package declares no asset dependency.  Editorial formatting only, in addition: an AMS \texttt{amsart} preamble that restates the article's own declarations and macros used by the retained lines, namely the environments \texttt{lemma} on separate counters, together with the standard packages those lines need.  The retained environments print in this document as Lemma 1 in order of appearance, whereas the article prints the same items with the numbering of its own full document.  The complete native source of the article as distributed in the arXiv e-print for this exact version is bundled unchanged as \texttt{sources/195414/source0.tex}.
\begin{lemma}\label{lem:stabilizer}
  Fix $F\in\mathcal{F}_K$.
  \begin{itemize}
    \item[$(1)$] The stabilizer $W_F$ is generated by all  simple reflections $r_k\in W_K$ that fix the barycenter of $F$.
    \item[$(2)$] If $P$ is non-degenerate, then for every $s \in W_K$, the intersection $s(F) \cap F$ is either $F$ or empty.
  \end{itemize}
\end{lemma}

\begin{lemma}\label{lem:equiv surjection} Let $G$ be a finite group, and $A,B$ be $\Q[G]$-modules. If a module homomorphism $f:A\to B$ is surjective, then the induced homomorphism between invariants $\bar{f}:A^G\to B^G$ is also surjective. 
\end{lemma} 

\begin{lemma}\label{lem:image}
  Recall that for $F\in\mathcal{F}_K$, we have
  $\Phi(X_{F,e})=\sum_{s\in W_K/W_F} X_{F,s}$,
  and similarly for $\phi(X_{F,e})$.
  \begin{itemize}
    \item[$(1)$] The algebra $\mathcal{SR}(P)^{W_K}$ is generated by $\Phi(X_{F,e})$ for all $F\in\mathcal{F}_K$.
    \item[$(2)$] The algebra $H^*(X_{P})^{W_K}$ is generated by $\phi(X_{F,e})$ for all $F\in\mathcal{F}_K$.
  \end{itemize}
\end{lemma}

\begin{proof}
  It suffices to show $(1)$, since the  quotient homomorphism $\mathcal{SR}(P)\to H^*(X_{P})$ induces a surjection $\mathcal{SR}(P)^{W_K}\to H^*(X_{P})^{W_K}$ by \cref{lem:equiv surjection}.

For a nonzero element in $\mathcal{SR}(P)^{W_K}$, we can choose its representative $f$ in 
$\Q[X_{F,s}: (F,s)\in\mathcal{F}]$ which is also $W_K$-invariant by \cref{lem:equiv surjection}. 
One can assume 
\begin{align*} 
  f= c\cdot X_{F_1,s_1}^{n_1}X_{F_2,s_2}^{n_2}\cdots X_{F_p,s_p}^{n_p}+\text{other terms $A$}, 
\end{align*} 
where $c\neq 0$ and $s_1(F_{1})\cap s_2(F_{2})\cap\cdots \cap s_p(F_{p})\ne \emptyset$. 
There exists  $\tilde{s}\in W_K$ such that 
$\tilde{s}\left(s_1(F_{1})\cap\cdots \cap s_p(F_{p})\right)=F_{1}\cap\cdots\cap F_{p}$. 
The stablizer of the face $F_{1}\cap\cdots\cap F_{p}$ is the parabolic subgroup $W':=W_{F_1}\cap\cdots\cap W_{F_p}$.
Using $W_K$-invariance, we may therefore assume 
\begin{align*} 
  f= c\cdot\biggl(\sum_{s\in {W_K}/{W'}} X_{F_1,s}^{n_1}X_{F_2,s}^{n_2}\cdots X_{F_p,s}^{n_p}\biggr)+A. 
\end{align*}
A non-empty face of the form $s(F_1\cap \cdots \cap F_p)$ can be written uniquely as
$t_1(F_{1})\cap\cdots \cap t_p(F_{p})$ with $t_i\in W_K/W_{F_i}$.
Hence
\begin{align*} 
  f&= c\cdot\biggl(\sum_{t_1,\ldots,t_{p}} X_{F_1,t_1}^{n_1}X_{F_2,t_2}^{n_2}\cdots X_{F_p,t_p}^{n_p}\biggr)+A+\mbox{other terms in $I_P$},\\ &=c\cdot\biggl(\sum_{t_1} X_{F_1,t_1}\biggr)^{n_1}\biggl(\sum_{t_2} X_{F_2,t_2}\biggr)^{n_2}\cdots \biggl(\sum_{t_p} X_{F_p,t_p}\biggr)^{n_p}+A+\mbox{other terms in $I_P$}, 
\end{align*} 
where $(t_1,\ldots,t_p)$ ranges over $\frac{W_K}{W_{F_1}}\times\cdots\times\frac{W_K}{W_{F_p}}$.
 The second equality follows from \cref{lem:stabilizer}.$(2)$ which implies that 
 $$X_{F_1,t_1}X_{F_1,w_1}\notin I\Leftrightarrow t_1(F_{1})\cap w_1(F_{1})\ne\emptyset\Leftrightarrow t_1=w_1\in {W_K}/{W_{F_1}}.$$ 
We may then proceed to the remaining term $A$ and argue by induction on the number of monomials.
This completes the proof.
\end{proof}

\end{document}
